\section{Geometry, observation and the complete Green defect}\label{sec:geometry}

The preceding proof uses a uniform relative approximation, for which
direct nonvanishing is the strongest immediate consumer. Green's formula
becomes useful when only an integrated boundary estimate is available.
We first separate the geometric facts from the arithmetic information
that must support them. All statements in this section concern one
fixed holomorphic field at one fixed time.

\subsection{A complete signed identity and a quadratic bound}

Let $D$ be a rectangle, $p\in D$, and let $F$ be analytic on a
neighborhood of $\overline D$, with $F(p)\ne0$. Normalize the positive
Green function by a logarithmic pole, so that the Green function of
the unit disk with pole zero is $-\log|z|$. Denote harmonic measure
at $p$ by $\omega_D^p$; it has total mass one.

\begin{lemma}[Harmonic normalization]\label{lem:normalization-green}
If $A$ is analytic and nonzero near $\overline D$, the Green defect
of $H=AF$ equals that of $F$. If there are no boundary zeros, this
defect is exactly
\begin{equation}\label{eq:green-identity}
 \int_{\partial D}\log|F|\dd\omega_D^p-\log|F(p)|
 =\sum_{z_j\in D}m_jG_D(p,z_j).
\end{equation}
\end{lemma}
\begin{proof}
$\log|A|$ is harmonic, so its boundary average equals its value at
$p$. Subtract from $\log|F|$ the finitely many interior logarithmic
poles, with their multiplicities, and solve the remaining harmonic
Dirichlet problem. The logarithmic-pole terms give precisely the
Green functions in~\eqref{eq:green-identity}.
\end{proof}

\begin{lemma}[Quadratic boundary consumer]\label{lem:quadratic-green}
Without excluding boundary zeros, one has
\begin{equation}\label{eq:quadratic-green}
 \sum_{z_j\in D}m_jG_D(p,z_j)
 \le\frac12\log\left(1+
 \frac{\displaystyle\int_{\partial D}|F-1|^2\dd\omega_D^p
                         -|F(p)-1|^2}{|F(p)|^2}\right).
\end{equation}
Consequently, if $|F(p)|\ge m>0$ and a target $K\subset D$ has
$\inf_KG_D(p,z)\ge g>0$, the condition
\begin{equation}\label{eq:quadratic-margin}
 \int_{\partial D}|F-1|^2\dd\omega_D^p<m^2(\ee^{2g}-1)
\end{equation}
excludes every zero in $K$.
\end{lemma}
\begin{proof}
The mean-value property gives $\int F\dd\omega_D^p=F(p)$.
On expanding squares, the linear terms cancel exactly, leaving
\[
 \int|F|^2\dd\omega_D^p
 =|F(p)|^2+\int|F-1|^2\dd\omega_D^p-|F(p)-1|^2.
\]
Map the disk conformally onto the rectangle with its origin at $p$.
On circles of radius less than one, Jensen's formula and concavity of
the logarithm bound the zero sum by half the logarithm of the boundary
mean of $|F|^2$, minus $\log|F(p)|$. As the radius increases to one,
the zero sums increase to the Green sum, while continuity of $|F|^2$
gives convergence of its boundary means. This proves
\eqref{eq:quadratic-green}, including possible boundary zeros.
A zero in $K$ would contribute at least $g$, contradicting the strict
bound in~\eqref{eq:quadratic-margin}.
\end{proof}

The subtraction of $|F(p)-1|^2$ is exact, not an estimated negative
term. The later application may discard it to obtain a simpler
sufficient inequality. A quadratic boundary bound can be much smaller
than a supremum bound. On the unit disk, for example,
$F_m(z)=1+2((1+z)/2)^m$ has supremum error two and boundary quadratic
error tending to zero by dominated convergence.

\subsection{Explicit contour and observation costs}

Write $k=a-\delta$, $b=a+\eta$, $d=b-k$, and take
\[
 D=\{|x-X|<W,\ \ell<y<u\},\qquad
 p=X+\ii b,\quad K=\{X+\ii y:k\le y\le a\}.
\]
Let $e_-=k-\ell>0$ and $e_+=u-b>0$.
Suppose for this paragraph that $|F-1|\le\varepsilon<1$ on
$\overline D$. Every point between $X+\ii k$ and $p$ has a Cauchy
disk of radius $r=\min(W,e_-,e_+)$. Thus
$|F'|\le\varepsilon/r$ and $|F(p)|\ge1-\varepsilon$.
The segment retains at least half that floor if
\begin{equation}\label{eq:cauchy-cost}
 d\le\frac{(1-\varepsilon)r}{2\varepsilon},
 \quad\text{or equivalently}\quad
 \varepsilon\le\frac{r}{2d+r}.
\end{equation}
At fixed vertical height $u-\ell>d$, equal vertical margins maximize
the available radius. This is a geometric optimization under a common
uniform error, not an acquisition of that error.

For a disk of radius $R>d/2$ centered halfway between $p$ and the
deepest target point, its Green function on the diameter gives the
worst target value
\begin{equation}\label{eq:disk-cost}
 g_{\rm disk}=\log\frac{R^2+d^2/4}{Rd}.
\end{equation}
Indeed the signed diameter coordinates of the two endpoints are
$d/2$ and $-d/2$, and the disk formula is
$\log|(R^2-st)/(R(s-t))|$.
It increases as the target moves upward toward the observation.
If $R=d/2+r$ fits inside the rectangle, domain monotonicity gives
$G_D\ge G_{\rm disk}$. The uniform-error logarithmic allowance is
$\log((1+\varepsilon)/(1-\varepsilon))$, so strict exclusion follows if
\begin{equation}\label{eq:disk-error-threshold}
 \varepsilon<\left(\frac{r}{d+r}\right)^2.
\end{equation}
This identity follows by solving the comparison of the two logarithms.
When a uniform error below one is already known on all of $D$, direct
nonvanishing excludes $D$ without either extra condition. Formulas
\eqref{eq:cauchy-cost}--\eqref{eq:disk-error-threshold} quantify the
costs of the alternative consumers and of future weaker inputs.

\subsection{The strip comparison and the cost of the ends}

Set $d_s=u-\ell$, $\alpha=\pi(b-\ell)/d_s$,
$\beta=\pi(y-\ell)/d_s$. Mapping the strip
$S=\{\ell<\Im z<u\}$ to the upper half-plane gives
\begin{equation}\label{eq:strip-green}
 G_S(p,x+\ii y)=\frac12\log
 \frac{\cosh(\pi(x-X)/d_s)-\cos(\alpha+\beta)}
      {\cosh(\pi(x-X)/d_s)-\cos(\alpha-\beta)}.
\end{equation}
On its centerline the minimum over $K$ is at $y=k$, with value
\[
 \log\frac{\sin(\pi(b+k-2\ell)/(2d_s))}
                {\sin(\pi(b-k)/(2d_s))}.
\]
For the rectangle with the same horizontal edges,
\begin{equation}\label{eq:strip-rectangle}
 0\le G_S(p,z)-G_D(p,z)
 \le\{cosh(\pi W/d_s)-1\}^{-1}.
\end{equation}
After cancelling the common pole, the difference is harmonic, vanishes
on the horizontal edges and is bounded on the vertical edges by the
right side, using $\log(1+q)\le q$. The maximum principle proves
the comparison. Enlarging a domain thus improves its Green signal,
but also changes the boundary data that must be paid.

The positive harmonic barrier
\[
 \frac{\cosh(\kappa_s(x-X))\cos(\kappa_s(y-y_m))}
 {\cosh(\kappa_s W)\cos(\kappa_s d_s/2)},
 \qquad y_m=(u+\ell)/2,\quad0<\kappa_s<\pi/d_s,
\]
is at least one on the two vertical sides and nonnegative on the
horizontal edges. Therefore their total harmonic measure at $p$ is
at most $\operatorname{sech}(\kappa_s W)/\cos(\kappa_s d_s/2)$.
The same barrier bounds the ends of a graph domain contained in this
strip. A passage to an infinite graph contour additionally requires
integrable boundary logarithms and a sequence of admissible end cuts.
For example, polynomial bounds on the absolute logarithms are absorbed
by this exponential end estimate. Geometry alone does not supply those
canonical bounds.

\subsection{Observation placement and a conditional scale rule}

With $D$ and $K$ fixed and the observation above $K$ on the centerline,
raising the observation decreases the worst Green signal. To see this,
map the symmetric domain to the disk so that its centerline maps
monotonically to a diameter. If $t(b)>t(y)$ are the real disk
coordinates, the pseudohyperbolic distance
$(t(b)-t(y))/(1-t(b)t(y))$ increases with $b$. Its negative logarithm
is the Green function. The Cauchy radius also decreases, while the
distance to the target increases. Under a common uniform error, the
smallest independently admissible observation buffer is therefore best.
With unequal boundary costs the choice can differ: the lower and upper
strip masses are $(u-b)/d_s$ and $(b-\ell)/d_s$.

A useful scale check is
\[
 k=3a/4,\quad b=9a/8,\quad\ell=a/2,\quad u=11a/8,
 \quad W\ge7a/16.
\]
Then $d=3a/8$, $r=a/4$, $R=7a/16$.
The Cauchy condition is $\varepsilon\le1/4$; the disk signal is
$\log(29/21)$ and its strict error threshold is $\varepsilon<4/25$.
The length and derivative costs cancel under rescaling. Conversely,
if a vertical margin is $o(d)$, these particular thresholds degenerate.
Even the exact Green signal is at most
\[
 \min\{\log(1+2e_-/d),\ \log(1+2e_+/d)\},
\]
by comparison with the two containing half-planes at the deepest point.

A canonical estimate with the same margin on every one of these
scaled collars would support $a_{n+1}=3a_n/4$. No such estimate is
proved here. If its onset were $X_n\to\infty$, only finitely many
steps would be justified at a fixed height. Both the carrier and the
arithmetic error must therefore be controlled uniformly enough for an
actual iteration.
